\begin{figure}[p]\centering\begin{tikzpicture}[gclplate]\node[anchor=west,font=\sffamily\bfseries\Large,gcllabel] at (-6,3.8){CUT REDUCTION};\draw[GCLWarm,line width=1pt](-6,3.45)--(6,3.45);\node[gcllabel,draw](h) at (0,2){cut rank $(r,h)$};\node[gcllabel,draw](l) at (0,.2){smaller rank $(r',h')$};\draw[gclbluearrow](h)--node[right,gcllabel]{principal/commuting step}(l);\end{tikzpicture}\caption{\textbf{Plate 21. Cut reduction trace.} A bounded cut step preserves the end sequent while decreasing the declared measure. Scope: computed trace is not the general cut-elimination theorem.}\label{plate:cut-trace}\end{figure}
